\clearpage\corrigesChapitre{14}{Émettre et transmettre un son}
\corrigesSouples
\begin{blocLivre}{cor-14-1}\begin{corrige}{1}{Une source sonore}
\begin{questions}
\item La corde se met en vibration.
\item La caisse aide à transmettre les vibrations à un plus grand volume d’air et peut renforcer le son émis.
\end{questions}
\end{corrige}\end{blocLivre}
\begin{blocLivre}{cor-14-2}\begin{corrige}{2}{Le rôle du milieu}
\schema[corrige]{propagation-son}
\begin{questions}
\item Le milieu matériel devient moins dense pour transmettre les vibrations ; d’autres transmissions par les supports peuvent subsister.
\item Non, le son nécessite un milieu matériel.
\end{questions}
\end{corrige}\end{blocLivre}
\begin{blocLivre}{cor-14-3}\begin{corrige}{3}{Une durée de propagation}
\begin{questions}
\item $\Delta t=d/v=68,0/340=\qty{0,200}{s}$.
\item $\qty{200}{ms}$.
\end{questions}
\end{corrige}\end{blocLivre}
\begin{blocLivre}{cor-14-4}\begin{corrige}{4}{Une distance sonore}
\begin{questions}
\item $d=v\Delta t=340\times0,0500=\qty{17,0}{m}$.
\item Non : $d=1500\times0,0500=\qty{75,0}{m}$.
\end{questions}
\end{corrige}\end{blocLivre}
\begin{blocLivre}{cor-14-5}\begin{corrige}{5}{Deux microphones}
\schema[corrige]{retard-son}
\begin{questions}
\item $v=1,02/(3,00\times10^{-3})=\qty{340}{m.s^{-1}}$.
\item Pour mesurer le retard de la même perturbation, sans confondre deux événements différents.
\end{questions}
\end{corrige}\end{blocLivre}
\begin{blocLivre}{cor-14-6}\begin{corrige}{6}{Un obstacle renvoie le son}
\schema[corrige]{echo}
\begin{questions}
\item $L=340\times0,300=\qty{102}{m}$.
\item $d=L/2=\qty{51,0}{m}$ car le signal effectue l’aller et le retour.
\end{questions}
\end{corrige}\end{blocLivre}
\begin{blocLivre}{cor-14-7}\begin{corrige}{7}{Comparer des vitesses}
\begin{questions}
\item $1500/340\simeq4,4$.
\item Le signal dans l’eau arrive plus tôt, si les émissions sont simultanées et les distances égales.
\end{questions}
\end{corrige}\end{blocLivre}
\begin{blocLivre}{cor-14-8}\begin{corrige}{8}{Une matière qui reste}
\begin{questions}
\item Non. Elles oscillent localement, sans déplacement global sur 10 m.
\item La perturbation mécanique et l’énergie associée.
\end{questions}
\end{corrige}\end{blocLivre}
\begin{blocLivre}{cor-14-9}\begin{corrige}{9}{Position de la source}
\begin{questions}
\item Non : le retard dépend de la différence des distances source-microphone, qui peut être plus petite.
\item Aligner source, premier microphone et second microphone, dans cet ordre.
\end{questions}
\end{corrige}\end{blocLivre}
\begin{blocLivre}{cor-14-10}\begin{corrige}{10}{Améliorer une résolution}
\begin{questions}
\item L’incertitude de lecture est importante devant le retard mesuré.
\item Augmenter la séparation dans une géométrie adaptée et limiter les échos parasites, puis répéter.
\end{questions}
\end{corrige}\end{blocLivre}
\begin{blocLivre}{cor-14-11}\begin{corrige}{11}{Un sonar}
\schema[corrige]{echo}
\begin{questions}
\item $d=1500\times0,0800/2=\qty{60,0}{m}$.
\item On trouverait 120 m, soit une distance deux fois trop grande.
\end{questions}
\end{corrige}\end{blocLivre}
\begin{blocLivre}{cor-14-12}\begin{corrige}{12}{Un retard mal lu}
\demarche{Mesurer une différence de dates, puis utiliser la distance supplémentaire parcourue.}
\begin{questions}
\item $\Delta t=15,0-12,0=\qty{3,0}{ms}$.
\item $v=1,02/0,0030=\qty{340}{m.s^{-1}}$. La date de 15,0 ms dépend de l’origine de l’acquisition ; seule la différence des dates décrit le trajet supplémentaire.
\end{questions}
\retiens{Une date d’acquisition dépend de l’origine du temps ; le retard de propagation est une différence.}
\end{corrige}\end{blocLivre}

